\documentclass[11pt]{article}

\usepackage[margin=1in]{geometry}
\usepackage{amsmath, amssymb, amsthm}
\usepackage{bm}
\usepackage{enumitem}
\usepackage{xcolor}
\usepackage[colorlinks=true, linkcolor=blue, urlcolor=blue]{hyperref}
\usepackage{parskip}

\newcommand{\ip}[2]{\langle #1,\, #2\rangle}
\newcommand{\db}{\Delta b}

\title{\vspace{-2.5em}Predicting fine-tuning generalization from activation geometry\\[0.4em]
\large Decoding Dan's $M$-update framing}
\author{}
\date{\today}

\begin{document}
\maketitle
\vspace{-2em}

\section*{The proposal}

Represent a \emph{context} as an activation vector $v_c$ (read at the token after the
last user-message token) and a \emph{behavior} as an activation vector $v_b$ (read at
the token after the assistant-message token). Suppose that, before training, the
model's context~$\to$~behavior map is approximately linear,
\[
  v_b' = M\,v_c'.
\]
Dan's conjecture: after fine-tuning context $c'$ to produce a new behavior $v_b''$
(instead of its original $v_b'$), the map becomes a rank-one correction of the old one,
\[
  \boxed{\;M_{\text{new}} \;=\; M \;+\; (v_b'' - v_b')\,v_c'^{\,\top}\;}
\]
The question is whether something of this shape predicts how training on
$c_0\!\to\!b_0$ generalizes to a different $c_1\!\to\!b_1$.

\section*{What the equation actually claims}

The update is a rank-one outer product: the behavior shift $(v_b''-v_b')$ times the
trained context $v_c'^{\,\top}$. Evaluate the new map at the training point:
\[
  M_{\text{new}}\,v_c'
  = M v_c' + (v_b''-v_b')\,\ip{v_c'}{v_c'}
  = v_b' + (v_b''-v_b')\,\|v_c'\|^2 .
\]
This exactly hits the target $v_b''$ only when $\|v_c'\|=1$, so the framing wants to
live in \emph{normalized / cosine} space with a fitted scalar gain. Now evaluate at a
\emph{new} context $c_1$:
\[
  v_b(c_1) = M v_{c_1} + (v_b''-v_b')\,\ip{v_c'}{v_{c_1}},
\]
so the predicted \textbf{change in behavior at any new context} is
\[
  \boxed{\;\Delta v_b(c_1) \;=\; (v_b''-v_b')\,\cdot\,\ip{v_c'}{v_{c_1}}\;}
\]

Two falsifiable claims fall out:

\begin{enumerate}[leftmargin=1.4em]
  \item \textbf{Direction is constant.} The behavior shifts in the \emph{same direction
  everywhere} --- the direction it was pushed at the training point. Only the magnitude
  varies across contexts.
  \item \textbf{Magnitude $\propto$ context cosine.} How much training on $c'$ bleeds
  into $c_1$ scales with the similarity $\ip{v_c'}{v_{c_1}}$ of their (base-model)
  context vectors.
\end{enumerate}

\section*{The clean bridge to what we already measure}

We do not measure $v_b$ directly; we measure $\Delta \log P(\text{marker})$ on-policy at
the post-response slot. If the behavior is linearly readable --- the persona-vector
premise --- there is a probe direction $w$ with $\ell(c) = \ip{w}{v_b(c)}$ for the
behavior logit. Then
\[
  \Delta \ell(c_1)
  = \ip{w}{\,v_b''-v_b'} \cdot \ip{v_c'}{v_{c_1}}
  = \underbrace{\ip{w}{\db}}_{\text{one scalar per behavior}} \cdot\; \ip{v_c'}{v_{c_1}}.
\]
That is a \textbf{one-free-parameter prediction}: leakage to every bystander context
should be \emph{linear in its base-model cosine to the trained context, with a single
shared slope}. Directly testable against the existing leakage tables (\#432$\to$\#456,
\#448): regress measured $\Delta\log P(\text{marker})$ per persona against base-model
context cosine; check linearity-through-origin with one slope.

\section*{This is the mechanism for \#404/\#458, not a new direction}

The persona-distance predictor line (\#404/\#458: does base-model KL/JS/cosine predict
leakage?) is the \emph{empirical} version of exactly this. Dan's $M$-update is the
generative model that says \emph{why} distance predicts leakage, and it sharpens the
existing result three ways:
\begin{itemize}[leftmargin=1.4em]
  \item predicts the relationship is \textbf{linear in cosine}, not merely monotonic;
  \item predicts a \textbf{shared slope} across contexts for a fixed behavior;
  \item adds genuinely new, untested content --- the \textbf{directional claim}: the
  activation-space shift is the \emph{same direction} everywhere. Current predictors
  only touch magnitude. Whether $\Delta v_b(c_i)$ is rank-one / constant-direction is the
  part nobody has checked.
\end{itemize}

\section*{It also explains contrastive negatives (\#18/\#207/\#383)}

Positive-only training gives $M + \db\,c_{\text{src}}^{\top}$. Leakage at bystander
$c_j$ is $\db\cdot\ip{c_{\text{src}}}{c_j}$; since persona contexts share a high baseline
cosine, you get \emph{uniform} leakage --- exactly the \#18/\#207 observation
(92--98\% on all bystanders). Contrastive negatives add \textbf{negative rank-one terms}
at the bystander contexts that cancel the $\ip{c_{\text{src}}}{c_j}$ term, which is why
the distance$\to$leakage gradient exists \emph{only} in the contrastive regime. The
framing predicts this for free:
\[
  M_{\text{new}} \approx M + \db\,c_{\text{src}}^{\top}
    - \sum_j \varepsilon_j\, c_{\text{byst},j}^{\top}.
\]

\section*{Where it is ``a bit wrong'' (Dan's hedge), ranked by impact}

\begin{enumerate}[leftmargin=1.4em]
  \item \textbf{Direction rotation is the real risk.} The strong claim is a
  context-independent shift direction. Plausibly it rotates --- evil-on-coding could
  land as ``evil'' nearby but ``off-target weird'' far away. First thing to measure:
  SVD the matrix of realized $\Delta v_b(c_i)$ and check it is rank-one with right vector
  $\approx v_c'$ and left vector $\approx \db$.
  \item \textbf{Rank-one sufficiency.} SGD distributes the update across all weights, so
  the induced map-change may be higher rank. ROME/MEMIT show narrow edits \emph{are}
  roughly rank-one, so it is a reasonable prior --- but it is the load-bearing
  assumption.
  \item \textbf{$M$ itself shifts $=$ emergent misalignment.} The model assumes $M$ is
  preserved except along $v_c'$. EM is precisely the failure of that: narrow training
  broadly rewrites $M$. In-framing, EM corresponds to either (a) a trained context
  $v_c'$ with high mean similarity to the whole context distribution (a generic context
  $\Rightarrow$ broad leakage), or (b) a realized shift with a large off-rank-one
  component. A concrete, testable account of \emph{when} narrow training generalizes
  broadly.
  \item \textbf{Normalization \& layer choice.} Needs cosine $+$ fitted gain, and a
  sweep over which layer / layer-pair gives the cleanest linear $M$. Context and behavior
  may read best at different layers.
  \item \textbf{Lossiness.} A single-token $v_c$/$v_b$ may not capture the full
  context/behavior, and emission goes through softmax $+$ sampling. The linear-probe
  bridge handles the last step, but only if the behavior is genuinely linearly readable.
\end{enumerate}

\section*{Extending to the non-linear case}

The before-map $v_b = M v_c$ is a linearization of the true forward map $f$ between the
two token positions. Two \emph{distinct} linearizations hide inside Dan's single $M$, and
separating them is what makes the non-linear extension non-trivial rather than cosmetic.

\paragraph{Input-space Jacobian (describes current behavior).} For a smooth $f$, locally
$f(v_c) \approx f(v_c') + J(v_c')(v_c - v_c')$ with $J = \partial v_b/\partial v_c$. So
the $M$ in Dan's \emph{before} equation is really the Jacobian of the frozen model at the
operating point. Its predictions hold only \emph{locally}; far from $v_c'$ the
second-order term dominates and the shift direction rotates --- exactly the
``direction-rotation'' failure mode, now derived rather than asserted.

\paragraph{Weight-space NTK (describes how training transfers).} The quantity Dan
actually wants --- how training on $c'$ generalizes to $c_1$ --- is not an input-space
object at all; it lives in \emph{weight} space. One gradient step on the regression loss
$\tfrac12\|f_\theta(c')-b''\|^2$ moves the function at any context $c$ by
\[
  \Delta f(c) = \eta\,(b'' - f(c'))\,
    \underbrace{\big\langle \nabla_\theta f(c),\ \nabla_\theta f(c')\big\rangle}
    _{\textstyle K(c,c')\ =\ \text{empirical NTK}},
\]
and the fully-trained (single-point, min-norm) solution is kernel regression:
\[
  \boxed{\;f_{\text{new}}(c) = f_{\text{old}}(c)
    + \frac{K(c,c')}{K(c',c')}\,\big(b'' - f_{\text{old}}(c')\big)\;}
\]

\paragraph{Dan's formula is exactly this for a linear model.} Take $f(c)=Mc$. Then
$\nabla_M f(c) = c$, so $K(c,c') = \ip{c}{c'}$ and, with $\|c'\|=1$,
\[
  f_{\text{new}}(c) = f_{\text{old}}(c) + \ip{c'}{c}\,(b'' - v_b'),
\]
which is Dan's $\Delta v_b(c) = (v_b''-v_b')\ip{v_c'}{v_c}$ verbatim. The dictionary:
\begin{itemize}[leftmargin=1.4em]
  \item behavior shift $(v_b''-v_b')$ $\;\leftrightarrow\;$ residual $b'' - f(c')$;
  \item context cosine $\ip{v_c'}{v_c}$ $\;\leftrightarrow\;$ normalized NTK
  $K(c,c')/K(c',c')$.
\end{itemize}
\textbf{So the non-linear extension is: replace the activation cosine with the empirical
NTK.} The raw cosine is the \emph{linear-model} NTK; the network's true NTK
$K(c,c')=\ip{\nabla_\theta f(c)}{\nabla_\theta f(c')}$ is the object that actually governs
generalization (Malladi et al., \emph{A Kernel-Based View of LM Fine-Tuning}, 2023). For
many training points, $f_{\text{new}}(c) = f_{\text{old}}(c) + \sum_i K(c,c_i)\,\alpha_i$
with $\alpha$ solving the kernel system --- the contrastive-negative cancellation falls
out as the off-diagonal of the kernel gram, \emph{derived} instead of assumed.

\paragraph{The deep payoff: lazy vs.\ rich.} The kernel picture is exact only in the
\emph{lazy} regime, where the NTK stays (approximately) constant through training. If
persona / marker fine-tuning is lazy, Dan's framing-with-NTK is fully predictive. If it
is \emph{rich} (feature-learning), training moves the kernel itself, and generalization
is \emph{not} fixed-kernel regression --- which is precisely the in-framing signature of
emergent misalignment ($K$ itself is rewritten; failure mode 3 above). So ``can we
predict generalization from base-model geometry?'' sharpens into the single diagnostic
\emph{``is persona fine-tuning in the lazy regime?''} --- whose answer decides whether
the whole program works.

\paragraph{Two cheaper intermediate extensions.}
\begin{itemize}[leftmargin=1.4em]
  \item \textbf{Link function (handles the softmax).} Keep the linear prediction in
  \emph{logit} space and apply the known output nonlinearity on top: emission rate
  $=\sigma\!\big(\ell_0 + \ip{w}{\db}\,\ip{v_c'}{v_c}\big)$. A GLM, not a new theory ---
  removes the softmax-lossiness objection for free.
  \item \textbf{Kernelize the activation similarity.} Stay in activation space but
  replace the raw cosine with $\ip{\phi(v_c')}{\phi(v_c)}$ for a feature map $\phi$ (a
  later / more-processed layer, a learned probe space, attention-pattern features). Let
  the data pick the similarity that best predicts leakage, without paying for the full
  NTK.
\end{itemize}

\paragraph{Testable.}
\begin{enumerate}[leftmargin=1.4em]
  \item Compute the empirical (LoRA-restricted, hence cheap) NTK gram between contexts;
  check whether $K(c',c_i)$ predicts held-out leakage \emph{better} than the raw cosine.
  The gap quantifies how much non-linearity matters.
  \item \textbf{Lazy-vs-rich diagnostic:} measure how much the NTK itself drifts between
  base and trained model across the persona panel. Large drift $\Rightarrow$ broad EM,
  and the fixed-kernel prediction should fail exactly where drift is largest.
\end{enumerate}

\paragraph{Caveats.} The empirical NTK is expensive in general (a parameter-gradient per
context); LoRA makes it tractable (gram over only the adapter parameters). The lazy
assumption is the whole ballgame --- but its \emph{failure} is informative, not fatal: it
localizes where feature-learning (hence broad generalization) actually happens.

\section*{Minimal experiment to nail it}

One pod, one clean test:
\begin{itemize}[leftmargin=1.4em]
  \item Fix a behavior (the marker is cleanest --- already measured on-policy). Train a
  single source context $c_0 \to$ marker.
  \item Across $N$ held-out contexts (the persona panel), measure: (a) realized
  $\Delta\log P(\text{marker})$; (b) realized activation shift $\Delta v_b(c_i)$;
  (c) \emph{base-model} context cosines $\ip{v_{c_0}}{v_{c_i}}$.
  \item \textbf{Headline test:} predict held-out leakage using \emph{only} base-model
  quantities plus the single training pair. If that prediction correlates with measured
  leakage across held-out contexts, the framing is predictive. Report (i) variance
  explained by the rank-one / cosine model, (ii) whether the shift direction is constant,
  (iii) the residual $=$ the off-target / direction-rotation part.
\end{itemize}
Even a partial result (``magnitude is cosine-predictable, direction rotates by
$X^\circ$'') is publishable and tells us something real.

\section*{Take}

Worth pursuing, as the \emph{deep} one. It (a) unifies three threads under one
generative model, (b) upgrades \#404/\#458 from ``distance correlates'' to ``here is the
mechanism $+$ a directional prediction,'' (c) is Dan's own framing, hence maximally
aligned, and (d) is exactly the ``fewer experiments, nailed carefully'' shape he keeps
asking for --- feeding the Aug 31 paper-draft milestone. Caveat if it stalls: do not
over-invest in the \emph{exact} rank-one form. The valuable deliverable is ``to what
extent is fine-tuning generalization linearly predictable from base-model activation
geometry''; the rank-one update is one hypothesis within that, not the whole prize.

\end{document}
